% 01-dat-van-de.tex ;  Sets up the empirical question
%
% Scope: brief recap of BPTT recurrence from ch01; the question this
% chapter answers by measurement, not derivation.

\section{Câu hỏi}
\label{sec:ch02-dat-van-de}

Trong BPTT, gradient của loss theo trạng thái ẩn tại bước $t$ được tính lùi
từ $T$ về $1$ bằng công thức truy hồi từ chương~01,
phương trình~\eqref{eq:bptt-dh-recurrence}:

\[
\pd{L}{\vect{h}_t}
= \mat{W}_{hy}\transpose \pd{L}{\vect{y}_t}
+ \mat{W}_{hh}\transpose \;
  \Bigl(
    \pd{L}{\vect{h}_{t+1}}
    \odot \bigl(\vect{1} - \vect{h}_{t+1}^2\bigr)
  \Bigr).
\]

Hệ quả trực tiếp: nếu $\|\mat{W}_{hh}\| < 1$ và các phần tử của
$\vect{1} - \vect{h}_t^2$ đều không vượt quá 1 (điều luôn đúng với
$\tanh$), thì $\|\partial L / \partial \vect{h}_t\|$ nhỏ dần sau mỗi
bước lùi. Sau $T$ bước, gradient truyền từ bước cuối về bước đầu bị
suy giảm theo hàm mũ của $T$.

Đó là lý thuyết. Nhưng lý thuyết không nói liệu điều này có thực sự xảy
ra trong một bài toán cụ thể hay không. Một RNN với $\mat{W}_{hh}$ được
khởi tạo ngẫu nhiên có thể có trị riêng lớn nhất vượt quá 1; một bài toán
có thể không cần gradient truyền xa đến vậy; và trong thực tế, mạng học
được những biểu diễn làm chậm quá trình suy giảm.

Chương này trả lời câu hỏi đó bằng thực nghiệm. Tôi sẽ huấn luyện RNN
trên hai bài toán được thiết kế để \emph{ép} mạng phải nhớ thông tin qua
nhiều bước, và tôi sẽ đo gradient ở từng bước để bạn thấy nó suy giảm
như thế nào, hoặc không.
